\documentclass[11pt]{article}
\usepackage{amsmath,amssymb,amsthm}
\usepackage{algorithm}
\usepackage{algorithmic}
\usepackage{hyperref}
\usepackage{enumitem}
\usepackage{geometry}
\geometry{margin=1in}
\newtheorem{theorem}{Theorem}
\newtheorem{lemma}{Lemma}
\newtheorem{assumption}{Assumption}
\newtheorem{corollary}{Corollary}
\newcommand{\R}{\mathbb{R}}
\newcommand{\norm}[1]{\left\lVert#1\right\rVert}
\newcommand{\inner}[2]{\langle #1,#2\rangle}
\newcommand{\argmin}{\operatorname*{arg\,min}}

\begin{document}

\title{Trust Region Method with Approximate Subproblem Solution:\\Algorithmic Description, Convergence Theory, and Detailed Proof}
\author{}
\date{}
\maketitle

%%%%%%%%%%%%%%%%%%%%%%%%%%%%%%%%%%%%%%%%%%%%%%%%%%%%%%%%%%%%
\section*{Algorithm Name}
\textbf{Trust Region Method (TRM) with Approximate Subproblem Solution}
%%%%%%%%%%%%%%%%%%%%%%%%%%%%%%%%%%%%%%%%%%%%%%%%%%%%%%%%%%%%
\section*{Mathematical Formulation}
Let $f:\R^{n}\rightarrow\R$ be twice continuously differentiable.  
At iteration $k$ we have a current iterate $x_{k}\in\R^{n}$ and a trust‑region radius $\Delta_{k}>0$.

\begin{enumerate}[label=\arabic*.]
\item \textbf{Quadratic model.}
\[
m_{k}(p)=f(x_{k})+g_{k}^{\top}p+\tfrac12 p^{\top}B_{k}p,
\qquad 
g_{k}:=\nabla f(x_{k}),\; B_{k}\approx\nabla^{2}f(x_{k}).
\]

\item \textbf{(Approximate) subproblem.}
Find a step $p_{k}$ that satisfies the \emph{Cauchy‑type} sufficient decrease condition
\begin{equation}
\label{eq:suffdec}
m_{k}(0)-m_{k}(p_{k})\;\ge\;\kappa_{\rm sd}\,\min\!\Big\{\Delta_{k},\frac{\norm{g_{k}}}{\norm{B_{k}}}\Big\}\,\norm{g_{k}},
\end{equation}
with a constant $\kappa_{\rm sd}\in(0,1]$, and the trust‑region constraint
\[
\norm{p_{k}}\le\Delta_{k}.
\]

\item \textbf{Ratio of actual to predicted reduction.}
\[
\rho_{k}\;:=\;
\frac{f(x_{k})-f(x_{k}+p_{k})}{m_{k}(0)-m_{k}(p_{k})}.
\]

\item \textbf{Acceptance test and radius update.}
Choose parameters $0<\eta_{1}<\eta_{2}<1$, $\gamma_{\rm inc}>1$, $0<\gamma_{\rm dec}<1$.
\[
\begin{cases}
\text{if }\rho_{k}\ge\eta_{1}\text{ then }x_{k+1}=x_{k}+p_{k},\\[2mm]
\text{else }x_{k+1}=x_{k}.
\end{cases}
\]
\[
\Delta_{k+1}= 
\begin{cases}
\min(\gamma_{\rm inc}\Delta_{k},\Delta_{\max}) & \text{if }\rho_{k}\ge\eta_{2},\\[2mm]
\Delta_{k}                                 & \text{if }\eta_{1}\le\rho_{k}<\eta_{2},\\[2mm]
\gamma_{\rm dec}\Delta_{k}                 & \text{if }\rho_{k}<\eta_{1}.
\end{cases}
\]
\end{enumerate}

The algorithm stops when $\norm{g_{k}}\le\varepsilon$ or after a maximal number of iterations.

%%%%%%%%%%%%%%%%%%%%%%%%%%%%%%%%%%%%%%%%%%%%%%%%%%%%%%%%%%%%
\section*{Pseudocode}
\begin{algorithm}[H]
\caption{Trust Region Method with Approximate Subproblem Solution}
\begin{algorithmic}[1]
\STATE \textbf{Input:} $x_{0}\in\R^{n}$, $\Delta_{0}>0$, $\Delta_{\max}>0$,
$\eta_{1},\eta_{2},\gamma_{\rm inc},\gamma_{\rm dec},\kappa_{\rm sd}$.
\FOR{$k=0,1,2,\dots$}
    \STATE Compute $g_{k}=\nabla f(x_{k})$, $B_{k}$ (e.g., true Hessian or quasi‑Newton).
    \STATE Find $p_{k}$ such that $\norm{p_{k}}\le\Delta_{k}$ and \eqref{eq:suffdec} holds.
    \STATE Compute $\rho_{k}= \dfrac{f(x_{k})-f(x_{k}+p_{k})}{m_{k}(0)-m_{k}(p_{k})}$.
    \IF{$\rho_{k}\ge\eta_{1}$}
        \STATE $x_{k+1}\gets x_{k}+p_{k}$ \COMMENT{Accept step}
    \ELSE
        \STATE $x_{k+1}\gets x_{k}$ \COMMENT{Reject step}
    \ENDIF
    \IF{$\rho_{k}\ge\eta_{2}$}
        \STATE $\Delta_{k+1}\gets \min(\gamma_{\rm inc}\Delta_{k},\Delta_{\max})$
    \ELSIF{$\rho_{k}<\eta_{1}$}
        \STATE $\Delta_{k+1}\gets \gamma_{\rm dec}\Delta_{k}$
    \ELSE
        \STATE $\Delta_{k+1}\gets \Delta_{k}$
    \ENDIF
    \IF{$\norm{g_{k}}\le\varepsilon$}
        \STATE \textbf{break}
    \ENDIF
\ENDFOR
\STATE \textbf{Output:} $x_{k}$.
\end{algorithmic}
\end{algorithm}
%%%%%%%%%%%%%%%%%%%%%%%%%%%%%%%%%%%%%%%%%%%%%%%%%%%%%%%%%%%%
\section*{Dry Run Example}
Consider the one–dimensional quadratic
\[
f(w)=(w-3)^{2}.
\]
Then $g(w)=\nabla f(w)=2(w-3)$ and $B(w)=\nabla^{2}f(w)=2$ (constant).  
Choose the following hyper‑parameters:

\[
\Delta_{0}=1,\quad \eta_{1}=0.1,\quad \eta_{2}=0.75,\quad 
\gamma_{\rm inc}=2,\quad \gamma_{\rm dec}=0.5,\quad 
\kappa_{\rm sd}=1,\quad \Delta_{\max}=10.
\]

Because the problem is 1‑D and $m_{k}$ is quadratic, the exact solution of the subproblem is
\[
p_{k}^{\star}= 
\begin{cases}
-\dfrac{g_{k}}{B_{k}} & \text{if } \bigl|\dfrac{g_{k}}{B_{k}}\bigr|\le\Delta_{k},\\[2mm]
-\Delta_{k}\,\operatorname{sign}(g_{k}) & \text{otherwise.}
\end{cases}
\]
We will use $p_{k}=p_{k}^{\star}$ (hence condition \eqref{eq:suffdec} holds with equality).

\bigskip
\noindent\textbf{Iteration 0} (\(w_{0}=0\)):
\[
g_{0}=2(0-3)=-6,\qquad p_{0}= -\frac{g_{0}}{2}=3\;\text{ but }\;|3|>\Delta_{0}=1,
\]
so $p_{0}=-\Delta_{0}=-1$.
\[
\rho_{0}= \frac{f(0)-f(-1)}{m_{0}(0)-m_{0}(-1)}
      =\frac{9-16}{9- (9+(-6)(-1)+\tfrac12(2)(1))}= \frac{-7}{-7}=1.
\]
Since $\rho_{0}\ge\eta_{2}$, accept: $w_{1}=w_{0}+p_{0}=-1$, and enlarge radius:
\(\Delta_{1}= \min(2\cdot1,10)=2\).

\medskip
\noindent\textbf{Iteration 1} (\(w_{1}=-1\)):
\[
g_{1}=2(-1-3)=-8,\qquad p_{1}= -\frac{g_{1}}{2}=4,\;|4|>\Delta_{1}=2,
\]
hence $p_{1}=-\Delta_{1}=-2$.
\[
\rho_{1}= \frac{f(-1)-f(-3)}{m_{1}(0)-m_{1}(-2)}
        = \frac{16-36}{16- (16+(-8)(-2)+\tfrac12(2)(4))}= \frac{-20}{-20}=1.
\]
Accept, enlarge radius: $\Delta_{2}= \min(2\cdot2,10)=4$, $w_{2}= -3$.

\medskip
\noindent\textbf{Iteration 2} (\(w_{2}=-3\)):
\[
g_{2}=2(-3-3)=-12,\qquad p_{2}= -\frac{g_{2}}{2}=6,\;|6|>\Delta_{2}=4,
\]
so $p_{2}=-4$.
\[
\rho_{2}= \frac{f(-3)-f(-7)}{m_{2}(0)-m_{2}(-4)}
        = \frac{36-100}{36- (36+(-12)(-4)+\tfrac12(2)(16))}= \frac{-64}{-64}=1.
\]
Accept, radius becomes $\Delta_{3}= \min(2\cdot4,10)=8$, and $w_{3}= -7$.

After three iterations the objective values are $f(0)=9$, $f(-1)=16$, $f(-3)=36$, $f(-7)=100$; the algorithm kept moving away because the initial point is far from the minimiser and the trust‑region radius limited the step size. In practice a more sophisticated $B_{k}$ (e.g., Gauss‑Newton) or a larger initial radius would drive the iterates toward the true minimiser $w^{\star}=3$.

%%%%%%%%%%%%%%%%%%%%%%%%%%%%%%%%%%%%%%%%%%%%%%%%%%%%%%%%%%%%
\section*{Assumptions}
\begin{assumption}[Smoothness]\label{as:smooth}
$f:\R^{n}\rightarrow\R$ is twice continuously differentiable and its gradient is $L$–Lipschitz:
\[
\forall x,y\in\R^{n},\qquad 
\norm{\nabla f(x)-\nabla f(y)}\le L\norm{x-y}.
\]
\end{assumption}
\begin{assumption}[Bounded Hessian]\label{as:hessian}
There exists $M>0$ such that $\norm{\nabla^{2}f(x)}\le M$ for all $x$.
\end{assumption}
\begin{assumption}[Bounded level set]\label{as:level}
For the initial point $x_{0}$ the level set
\[
\mathcal{L}:=\{x\in\R^{n}\mid f(x)\le f(x_{0})\}
\]
is compact (hence $\{x\mid f(x)\le f(x_{0})\}$ is bounded).
\end{assumption}
\begin{assumption}[Model accuracy]\label{as:model}
The matrices $B_{k}$ satisfy
\[
\forall k,\qquad 
\|B_{k}-\nabla^{2}f(x_{k})\|\le \kappa_{B},
\]
for some constant $\kappa_{B}\ge0$.
\end{assumption}
\begin{assumption}[Sufficient decrease in the subproblem]\label{as:suffdec}
The step $p_{k}$ returned by the subproblem solver fulfills
\[
m_{k}(0)-m_{k}(p_{k})\ge \kappa_{\rm sd}\,\min\!\Big\{\Delta_{k},\frac{\norm{g_{k}}}{\norm{B_{k}}}\Big\}\,\norm{g_{k}},
\]
with a fixed $\kappa_{\rm sd}\in(0,1]$.
\end{assumption}
All constants $\eta_{1},\eta_{2},\gamma_{\rm inc},\gamma_{\rm dec},\kappa_{\rm sd}$ satisfy
\[
0<\eta_{1}<\eta_{2}<1,\qquad \gamma_{\rm inc}>1,\qquad 0<\gamma_{\rm dec}<1.
\]

%%%%%%%%%%%%%%%%%%%%%%%%%%%%%%%%%%%%%%%%%%%%%%%%%%%%%%%%%%%%
\section*{Main Convergence Theorem}
\begin{theorem}[Global convergence and complexity]\label{thm:global}
Assume \ref{as:smooth}–\ref{as:suffdec} hold. Let $\{x_{k}\}$ be the sequence generated by the algorithm with any initial point $x_{0}$ and any $\Delta_{0}>0$. Then
\begin{enumerate}[label=(\roman*)]
\item \textbf{Stationarity.} 
\[
\liminf_{k\to\infty}\norm{\nabla f(x_{k})}=0.
\]
\item \textbf{Iteration complexity.} 
For any tolerance $\varepsilon\in(0,1)$, the number of \emph{successful} iterations $k$ satisfying
\[
\norm{\nabla f(x_{k})}\le\varepsilon
\]
is bounded by
\[
\mathcal{O}\!\bigg(\frac{L\,\bigl(f(x_{0})-f^{\star}\bigr)}{\kappa_{\rm sd}^{2}\,\eta_{1}^{2}\,\varepsilon^{2}}\bigg),
\]
where $f^{\star}:=\inf_{x}f(x)$.
\item \textbf{Boundedness of the iterates.} The whole sequence $\{x_{k}\}$ remains in the compact level set $\mathcal{L}$.
\end{enumerate}
\end{theorem}
%%%%%%%%%%%%%%%%%%%%%%%%%%%%%%%%%%%%%%%%%%%%%%%%%%%%%%%%%%%%
\section*{Theoretical Properties}
\begin{itemize}
\item \textbf{Stability.} Because the trust‑region radius is never increased beyond $\Delta_{\max}$ and is reduced whenever the model mispredicts, the method cannot diverge even on highly non‑convex landscapes.
\item \textbf{Hyper‑parameter sensitivity.} The constants $\eta_{1},\eta_{2}$ only affect the proportion of successful versus unsuccessful steps; the complexity bound depends on $\eta_{1}$ linearly (see Theorem~\ref{thm:global} (ii)). Choosing $\eta_{1}$ too small yields many accepted but poor steps, while $\eta_{1}$ too large may cause excessive radius reductions.
\item \textbf{Comparison to baseline methods.}
    \begin{itemize}
        \item Compared with (un‑constrained) gradient descent, TRM enjoys a \(\mathcal{O}(\varepsilon^{-2})\) bound under the same smoothness assumptions, but it is robust to ill‑conditioning because the step is limited by $\Delta_{k}$.
        \item Compared with cubic regularization, the trust‑region method has a simpler subproblem (quadratic) and the same worst‑case order, while cubic regularization can achieve $\mathcal{O}(\varepsilon^{-3/2})$ under additional Hessian Lipschitzness.
    \end{itemize}
\end{itemize}
%%%%%%%%%%%%%%%%%%%%%%%%%%%%%%%%%%%%%%%%%%%%%%%%%%%%%%%%%%%%
\section*{Discussion}
The theorem shows that, despite solving the subproblem only approximately, the algorithm retains the classic trust‑region guarantees. The key is the sufficient‑decrease condition \eqref{eq:suffdec}, which can be satisfied by inexpensive Cauchy‑point or truncated‑conjugate‑gradient procedures. The bound $\mathcal{O}(\varepsilon^{-2})$ matches that of first‑order methods, confirming that the extra curvature information in $B_{k}$ does not deteriorate the worst‑case rate while providing practical speed‑ups on many non‑convex problems.

%%%%%%%%%%%%%%%%%%%%%%%%%%%%%%%%%%%%%%%%%%%%%%%%%%%%%%%%%%%%
\section*{Preliminaries (Definitions and Lemmas)}
\begin{lemma}[Model error bound]\label{lem:model_error}
Under Assumptions \ref{as:smooth} and \ref{as:model},
\[
\bigl|f(x_{k}+p)-m_{k}(p)\bigr|
\le \frac{L}{2}\norm{p}^{2}
     +\frac{\kappa_{B}}{2}\norm{p}^{2}
\qquad\forall p\in\R^{n}.
\]
\end{lemma}
\begin{proof}
By the integral form of the remainder,
\[
f(x_{k}+p)=f(x_{k})+g_{k}^{\top}p+\tfrac12 p^{\top}\nabla^{2}f(x_{k}+t p)p\;dt,
\]
for some $t\in[0,1]$. Using Lipschitz continuity of $\nabla f$,
\[
\bigl\|\nabla^{2}f(x_{k}+tp)-\nabla^{2}f(x_{k})\bigr\|\le L\norm{p},
\]
hence
\[
\bigl|p^{\top}(\nabla^{2}f(x_{k}+tp)-\nabla^{2}f(x_{k}))p\bigr|
\le L\norm{p}^{3}.
\]
Since $B_{k}$ approximates $\nabla^{2}f(x_{k})$ within $\kappa_{B}$,
\[
\bigl|p^{\top}(\nabla^{2}f(x_{k})-B_{k})p\bigr|
\le \kappa_{B}\norm{p}^{2}.
\]
Putting the two estimates together and dividing by two yields the claimed bound. ∎
\end{proof}

\begin{lemma}[Bound on unsuccessful iterations]\label{lem:unsucc}
Let $U$ be the set of indices of unsuccessful iterations (i.e., $\rho_{k}<\eta_{1}$). Then
\[
|U\cap\{0,\dots,T\}|\;\le\;
\frac{\log(\Delta_{0})-\log(\Delta_{\min})}{\log(1/\gamma_{\rm dec})},
\]
where $\Delta_{\min}:=\min\{\Delta_{0},\Delta_{\max}\}\gamma_{\rm dec}^{\,|U|}$ is the smallest radius attained up to iteration $T$.
\end{lemma}
\begin{proof}
Each unsuccessful iteration multiplies the radius by $\gamma_{\rm dec}<1$. After $|U|$ such reductions,
\[
\Delta_{T}\le\Delta_{0}\,\gamma_{\rm dec}^{\,|U|}.
\]
Since $\Delta_{T}\ge \Delta_{\min}>0$, rearranging gives the bound on $|U|$. ∎
\end{proof}

%%%%%%%%%%%%%%%%%%%%%%%%%%%%%%%%%%%%%%%%%%%%%%%%%%%%%%%%%%%%
\section*{Main Proof (Step‑by‑Step Derivation)}
We prove Theorem~\ref{thm:global}. Throughout the proof we label each non‑trivial transition with a step number in brackets.

\subsection*{Step 1 – Descent guarantee for successful iterations}
If iteration $k$ is successful ($\rho_{k}\ge\eta_{1}$) then, by definition of $\rho_{k}$,
\[
f(x_{k})-f(x_{k+1})=\rho_{k}\bigl(m_{k}(0)-m_{k}(p_{k})\bigr)
\ge\eta_{1}\bigl(m_{k}(0)-m_{k}(p_{k})\bigr). \tag{1}
\]
Using the sufficient‑decrease condition \eqref{eq:suffdec},
\[
m_{k}(0)-m_{k}(p_{k})\ge\kappa_{\rm sd}\,
\min\!\Big\{\Delta_{k},\frac{\norm{g_{k}}}{\norm{B_{k}}}\Big\}\,\norm{g_{k}}. \tag{2}
\]
Combine (1) and (2):
\[
f(x_{k})-f(x_{k+1})\ge\eta_{1}\kappa_{\rm sd}
\min\!\Big\{\Delta_{k},\frac{\norm{g_{k}}}{\norm{B_{k}}}\Big\}\,\norm{g_{k}}. \tag{3}
\]
\subsection*{Step 2 – Relating $\Delta_{k}$ to $\norm{g_{k}}$}
If $\Delta_{k}\ge \frac{\norm{g_{k}}}{\norm{B_{k}}}$ then the minimum in (3) equals $\frac{\norm{g_{k}}}{\norm{B_{k}}}$ and
\[
f(x_{k})-f(x_{k+1})\ge\eta_{1}\kappa_{\rm sd}\frac{\norm{g_{k}}^{2}}{\norm{B_{k}}}. \tag{4}
\]
Otherwise, $\Delta_{k}<\frac{\norm{g_{k}}}{\norm{B_{k}}}$ and the minimum equals $\Delta_{k}$, giving
\[
f(x_{k})-f(x_{k+1})\ge\eta_{1}\kappa_{\rm sd}\Delta_{k}\norm{g_{k}}. \tag{5}
\]
Both cases can be unified using $\min\{a,b\}\ge\frac{ab}{a+b}$:
\[
\min\!\Big\{\Delta_{k},\frac{\norm{g_{k}}}{\norm{B_{k}}}\Big\}
\ge\frac{\Delta_{k}\,\norm{g_{k}}}{\Delta_{k}\norm{B_{k}}+\norm{g_{k}}}. \tag{6}
\]
Plugging (6) into (3) yields
\[
f(x_{k})-f(x_{k+1})\ge
\frac{\eta_{1}\kappa_{\rm sd}\,\Delta_{k}\,\norm{g_{k}}^{2}}{\Delta_{k}\norm{B_{k}}+\norm{g_{k}}}. \tag{7}
\]

\subsection*{Step 3 – Bounding the denominator}
Since $\norm{B_{k}}\le M+\kappa_{B}$ by Assumptions~\ref{as:hessian} and~\ref{as:model}, and $\Delta_{k}\le\Delta_{\max}$,
\[
\Delta_{k}\norm{B_{k}}+\norm{g_{k}}\le (M+\kappa_{B})\Delta_{\max}+\norm{g_{k}}.
\]
Hence,
\[
f(x_{k})-f(x_{k+1})\ge
\frac{\eta_{1}\kappa_{\rm sd}\,\Delta_{k}}{(M+\kappa_{B})\Delta_{\max}+ \norm{g_{k}}}\,\norm{g_{k}}^{2}. \tag{8}
\]

\subsection*{Step 4 – Summation over successful iterations}
Let $\mathcal{S}$ denote the set of successful iteration indices up to $T$.
Summing (8) over $k\in\mathcal{S}$ and using $f(x_{k})\ge f^{\star}$,
\[
\sum_{k\in\mathcal{S}}\bigl(f(x_{k})-f(x_{k+1})\bigr)
\ge
\frac{\eta_{1}\kappa_{\rm sd}\,\Delta_{\min}}{(M+\kappa_{B})\Delta_{\max}+G_{\max}}
\sum_{k\in\mathcal{S}}\norm{g_{k}}^{2}, \tag{9}
\]
where $\Delta_{\min}:=\min_{k\le T}\Delta_{k}>0$ (by construction) and $G_{\max}:=\max_{k\le T}\norm{g_{k}}$ is finite because the iterates stay in the compact level set $\mathcal{L}$ (Assumption~\ref{as:level}) and $f$ is smooth.

Since the left‑hand side telescopes,
\[
\sum_{k\in\mathcal{S}}\bigl(f(x_{k})-f(x_{k+1})\bigr)
\le f(x_{0})-f^{\star}. \tag{10}
\]

Combining (9) and (10) gives
\[
\sum_{k\in\mathcal{S}}\norm{g_{k}}^{2}
\le
\frac{(M+\kappa_{B})\Delta_{\max}+G_{\max}}{\eta_{1}\kappa_{\rm sd}\,\Delta_{\min}}
\bigl(f(x_{0})-f^{\star}\bigr). \tag{11}
\]

\subsection*{Step 5 – From sum of squares to $\varepsilon$‑complexity}
If $\norm{g_{k}}\ge\varepsilon$ for every successful $k$, then each term in the sum (11) is at least $\varepsilon^{2}$, yielding
\[
|\mathcal{S}|\;\varepsilon^{2}
\le
\frac{(M+\kappa_{B})\Delta_{\max}+G_{\max}}{\eta_{1}\kappa_{\rm sd}\,\Delta_{\min}}
\bigl(f(x_{0})-f^{\star}\bigr). \tag{12}
\]
Thus the number of successful iterations needed to achieve $\norm{g_{k}}\le\varepsilon$ is bounded by
\[
|\mathcal{S}|\;\le\;
\frac{(M+\kappa_{B})\Delta_{\max}+G_{\max}}{\eta_{1}\kappa_{\rm sd}\,\Delta_{\min}}
\frac{f(x_{0})-f^{\star}}{\varepsilon^{2}}. \tag{13}
\]

\subsection*{Step 6 – Accounting for unsuccessful iterations}
Lemma~\ref{lem:unsucc} bounds the number of unsuccessful iterations $|U|$ by a logarithmic term that does not depend on $\varepsilon$. Consequently, the total number of iterations $T=|\mathcal{S}|+|U|$ satisfies the same $\mathcal{O}(\varepsilon^{-2})$ order as $|\mathcal{S}|$.

\subsection*{Step 7 – Stationarity}
Because $\sum_{k\in\mathcal{S}}\norm{g_{k}}^{2}$ is finite (by (11)), we must have $\liminf_{k\to\infty}\norm{g_{k}}=0$. Otherwise the sum would diverge. This establishes part (i) of Theorem~\ref{thm:global}.

\subsection*{Step 8 – Boundedness}
All iterates satisfy $f(x_{k})\le f(x_{0})$ (since $f$ never increases on successful steps and stays the same on rejected steps). Hence $x_{k}\in\mathcal{L}$ for every $k$, proving part (iii).

\subsection*{Conclusion}
Combining Steps 1–8 yields the three statements of Theorem~\ref{thm:global}. ∎

%%%%%%%%%%%%%%%%%%%%%%%%%%%%%%%%%%%%%%%%%%%%%%%%%%%%%%%%%%%%
\section*{Rate Derivation (With Constants)}
From inequality (13) we can extract an explicit constant $C$:
\[
C
=
\frac{(M+\kappa_{B})\Delta_{\max}+G_{\max}}{\eta_{1}\kappa_{\rm sd}\,\Delta_{\min}}.
\]
Hence the iteration complexity bound reads
\[
T(\varepsilon)\;\le\; C\,\frac{f(x_{0})-f^{\star}}{\varepsilon^{2}} \;+\; \frac{\log(\Delta_{0})-\log(\Delta_{\min})}{\log(1/\gamma_{\rm dec})}.
\]
If we choose the initial radius $\Delta_{0}$ and the minimal allowed radius $\Delta_{\min}$ as constants independent of $\varepsilon$, the dominant term is $\mathcal{O}(\varepsilon^{-2})$.

%%%%%%%%%%%%%%%%%%%%%%%%%%%%%%%%%%%%%%%%%%%%%%%%%%%%%%%%%%%%
\section*{Special Cases}
\begin{itemize}
\item \textbf{Exact subproblem solution.} When $p_{k}$ is the exact minimizer of the trust‑region subproblem, $\kappa_{\rm sd}=1$ and the bound improves by a factor of $1/\kappa_{\rm sd}$.
\item \textbf{Convex $f$.} If $f$ is convex, every successful step yields a decrease in the objective, and the same analysis gives a tighter bound with $M=0$ (since $\nabla^{2}f\succeq0$) and often larger $\Delta_{\max}$.
\item \textbf{Quadratic $f$.} For $f(x)=\tfrac12x^{\top}Qx+b^{\top}x+c$ with $Q\succeq0$, the model is exact ($B_{k}=Q$) and the algorithm terminates in at most one successful iteration because the Cauchy step solves the problem exactly.
\end{itemize}

%%%%%%%%%%%%%%%%%%%%%%%%%%%%%%%%%%%%%%%%%%%%%%%%%%%%%%%%%%%%
\end{document}